\scanpage{151}
On dit alors que $\Delta$ (ou la coalgèbre \inserted{$M$}) est anticommutative si \deleted{\unsure{tt}} la multiplication sur \deleted{$V(M)$} l'algèbre graduée
\[
\deleted{V(M)=V(M)^{\mathrm{pair}}}\qquad
V(M)=V(M)^{\mathrm{pair}}\times V(M)^{\mathrm{impair}}
\]
l'est, i.e. si on a (pour $x,y\in V(M)(A)$ hom mod 2)
\[
xy=(-1)^{\deg x\deg y}yx
\]
(i.e. $xy=yx$ si $x$ \unsure{ou|et} $y$ \deleted{ont $=$ degré mod 2}\inserted{sont de degré pair}, \deleted{$xy=yx$} $xy=-yx$ i.e. $xy+yx=0$ si l'un et l'autre de degré impair). On laisse au lecteur le soin de diagrammatiser cette condition.

On dit que $\Delta$ (ou la coalgèbre $M$) est alternée si la multiplication de $V(M)$ l'est, i.e. si on a pour \deleted{\unsure{alg $V$}} $x,y\in V(M)(A)$ homogènes\deleted{\unsure{et}} mod 2, on a
\[
\begin{cases}
xy=yx&\text{si }x\text{ ou }y\text{ de degré pair}\\
\textcircled{$x^2=0$}&\text{si }x\text{ de degré impair}
\end{cases}
\]
\margin{signifie que le composé
\[
\deleted{\unsure{M^{\mathrm{imp}}\longrightarrow1}}
\]
\diagram{\begin{tikzcd}
\unsure{M^0}\arrow[r,"\unsure{\Delta^{00}|\Delta^{11}}"] & M^1\otimes M^1\arrow[r] & \operatorname{Sym}^2M^1
\end{tikzcd}}
est nul (degrés mod 2)
}
(alors, si $x$ et $y$ de degré impair, utilisant $x^2=y^2=(x+y)^2=0$, on trouve $xy+yx=0$, donc l'Algèbre est anticommu., donc \inserted{pour} la coalgèbre aussi
\[
\text{alternée}\ \Longrightarrow\ \text{anticommutative}.
\]
Bien sûr, anticommutativité et \inserted{le caractère} alterné \deleted{\unsure{procèdent}} des \deleted{\unsure{tt fait}} coalgèbres se conservent par \uline{lim} quelc.

On définit les $k$-\inserted{co}dérivations d'une coalgèbre, et les $k$-antidérivations (tjrs
\margin{150}
